\section{Useful Conclusions and Fundamental Definitions}

\subsection{Probability Theory}
\begin{theorem}[Kolmogorove's Theorem] The probability distribution functions
$(F_{\textbf{t}})_{\textbf{t}\in\mathcal{T}}$ are the distribution functions of some stochastic process if and
only if for any integer $n$ and $\textbf{t} = (t_1 , \cdots , t_n) \in \mathcal{T}$ and $1\leq i\leq n$,
\[\lim_{x_i \to\infty} F_{\textbf{t}}(x) = F_{\textbf{t}(i)}(x(i))\]
where $\textbf{t}(i)$ and $x(i)$ are the $(n- 1)$-component vectors obtained by deleting the
$i$ th components of $\textbf{t}$ and $x$ respectively. 
\end{theorem}

\begin{theorem}[Spectral Theorem for Compact Self-Adjoint Operators]
Let $H$ be a real or complex Hilbert space, and let
\(
T:H\to H
\)
be a compact self-adjoint operator. Then there exists an orthonormal family
$\{e_n\}$ consisting of eigenvectors of $T$ corresponding to nonzero
eigenvalues $\{\lambda_n\}$ such that
\[
H=\ker T\oplus \overline{\operatorname{span}\{e_n\}}.
\]
Moreover, the following properties hold:
\begin{enumerate}
    \item Every eigenvalue $\lambda_n$ is real.
    \item Every nonzero eigenspace is finite-dimensional.
    \item The set of nonzero eigenvalues is at most countable, and if there
    are infinitely many of them, then $0$ is their only possible accumulation
    point.
    \item For every $x\in H$,
    \[
    Tx=\sum_n \lambda_n \langle x,e_n\rangle e_n.
    \]
\end{enumerate}
In particular, if $\ker T=\{0\}$, then $\{e_n\}$ is an orthonormal basis of
$H$.
\end{theorem}
\begin{proof}
    See Appendix $\eqref{prove_specT}$.
\end{proof}


\subsection{Definitions}

\begin{definition}
    [Stationarity] The time series $(X_t)_{t\in\mathbb{Z}}$ is said to be stationary if
    \begin{itemize}
        \item $\mathbb{E}|X_t|^2 < \infty$ for all $t\in\mathbb{Z}$.
        \item $\mathbb{E}X_t = m$ for all $t\in\mathbb{Z}$.
        \item $\text{Cov}(X_s,X_t)=:\gamma_X(s,t) = \gamma_X(s+h,r+h)$ for all $s,t,h\in\mathbb{Z}$.    
        \end{itemize}
\end{definition}

\begin{definition}
    [Strict Stationarity] The time series $(X_t)_{t\in\mathbb{Z}}$ is said to be
strictly stationary if the joint distributions of $(X_{t_1}, \cdots , X_{t_k})$ and $(X_{t_1+h}, \cdots , X_{t_k+h})$
are the same for all positive integers $k$ and for all $t_1, \cdots , t_k, h \in \mathbb{Z}$. 
\end{definition}

The strict stationarity may imply the stationarity easily.

\begin{definition}
    [Elementary Properties]. For a stationary process, we define the autocovariance $\gamma_X(h) = \gamma_X(t+h,t)$ for any $t\in\mathbb{Z}$ which is well-defined.\par If $\gamma$ is the autocovariance function of a stationary process $(X_t)_{t\in\mathbb{Z}}$, then
    \[
    \gamma(0) \geq 0, \quad |\gamma(h)| \leq \gamma(0)\text{ for all }h\in\mathbb{Z}
    \]
    and $\gamma(h) =  \gamma(-h)$ for any $h\in\mathbb{Z}$.
\end{definition}
\begin{proof}
    We only prove that $\gamma(0) \geq |\gamma(h)|$ for any $h\in\mathbb{Z}$.
    \[
    \begin{aligned}
        \gamma(h)^2 &= \left(\mathbb{E}(X_0-\mu(0))(X_{h} - \mu(h))\right)^2 \\
        &\leq \mathbb{E}(X_0-\mu(0))^2\mathbb{E}(X_h-\mu(h))^2 = \gamma(0)^2
    \end{aligned}
    \]
    by Cauchy-Schwarz Inequality.
\end{proof}

\begin{definition}
    [Non-Negative Definiteness] A real-valued function on the
integers, $\kappa : \mathbb{Z} \to \mathbb{R}$ is said to be non-negative definite if and only if 
\[\sum\limits_{i,j=1}^n a_i\kappa(t_i,t_j)a_j \geq 0\]
for all positive integers $n$ and for all vectors $a = (a_1 , \cdots , a_n) \in \mathbb{R}^n$ and $t =
(t_1, \cdots , t_n) \in \mathbb{Z}$. 
\end{definition}
\textbf{Remark}. Replace $t_i$ with $i$ in the definition above will obtain an equivalent conclusion.

\begin{theorem}
     [Characterization of Autocovariance Functions] A real-valued
function defined on the integers is the autocovariance function of a
stationary time series if and only if it is even and non-negative definite. 
\end{theorem}

\begin{definition}
    [Sample Autocovariance Function] The sample autocovariance function of $\{x_ 1 , \cdots , x_n\}$ is defined by
    \[
    \hat{\gamma}(h) := \dfrac{1}{n}\sum\limits_{i=1}^{n-h}(x_i-\bar{x})(x_{i+h}-\bar{x}), 0\leq h <n
    \] 
and $\gamma(h) = \gamma( -h), -n < h \leq 0$.
\end{definition}

\begin{proposition}
    The covariance matrix is symmetric and non-negative definite for any random vector.
\end{proposition}

\begin{proposition}
    If $\textbf{a}$ is an $m$-component column vector, $B$ is an $m \times n$ matrix, and $\textbf{X} =(X_1 , \cdots , X_n)$ where $\mathbb{E}|X_i|^2 < \infty$, i = $1, \cdots , n$, then the random vector,
    \[
    \textbf{Y} = \textbf{a} + B\textbf{X}
    \]
    has the mean $\mathbb{E}\textbf{Y} = \textbf{a} + B\mathbb{E}\textbf{X}$ and covariance matrix
    \[
    \Sigma_{\textbf{Y}\textbf{Y}} = B\Sigma_{\textbf{X}\textbf{X}}B^T
    \]
\end{proposition}

\begin{definition}[Multivariate Normal Distribution]   The random vector
$\textbf{Y} = ( Y_1 , \cdots , Y_n)$ is said to be multivariate normal, or to have a multivariate normal distribution, if and only if there exist a column vector $\textbf{a}$, a matrix $B$
and a random vector $\textbf{X}= (X_1 ,\cdots,X_m)$ with independent standard normal components, such that $\textbf{Y} = \textbf{a}+B\textbf{X}$, with
\[
\mathbb{E}\textbf{Y} = \textbf{a},\quad \Sigma\textbf{Y} = BB^T 
\]
\end{definition}

\begin{proposition}
    If $\textbf{Y} = (Y_1 , \cdots , Y_n)$ is a multivariate normal random vector such that $\mathbb{E}\textbf{Y} = \mu$ and $\Sigma_{\textbf{Y}\textbf{Y}} = \Sigma$, then the characteristic function of $\textbf{Y}$ is
    \[
    \phi_{\textbf{Y}}(\textbf{x}) = \exp\left(i\textbf{x}^T\mu - \dfrac{1}{2}\textbf{x}^T\Sigma\textbf{x}\right)
    \]
    If $\det \Sigma > 0$ then $\textbf{Y}$ has the density
    \[
f_{\mathbf{Y}}(\mathbf{y})
=
(2\pi)^{-n/2}
(\det{\Sigma})^{-1/2}
\exp\left[
-\frac{1}{2}
(\mathbf{y}-{\mu})^\prime
{\Sigma}^{-1}
(\mathbf{y}-{\mu})
\right].
\tag{1.6.12}
\]
\end{proposition}
\begin{proof}
    Consider $\textbf{X}$ to be an independent standard normal vector, then we know $\mu + \Sigma^{1/2}\textbf{X} \overset{d}{=} \textbf{Y}$, then
    \[
    \begin{aligned}
        \phi_{\textbf{Y}}(\textbf{x}) &= \mathbb{E} e^{i\textbf{x}^T\textbf{Y}} 
        = \mathbb{E} e^{i\textbf{x}^T(\mu + \Sigma^{1/2}\textbf{X})} \\
        &= \exp(i\textbf{x}^T\mu)\phi_{\textbf{X}}(\Sigma^{1/2}\textbf{x}) = \exp\left(i\textbf{x}^T\mu - \dfrac{1}{2}\textbf{x}^T\Sigma\textbf{x}\right)
    \end{aligned}
    \]
    
\end{proof}